\begin{figure}[htbp]
\centering
\begin{tikzpicture}[x=1cm,y=1cm,
  mk/.style={circle,fill=draftblue,text=white,font=\scriptsize\bfseries,inner sep=0pt,minimum size=4.2mm},
  lay/.style={font=\scriptsize,align=center}]
% geometry (schematic, not to scale)
\def\W{13.0}\def\Lx{4.0}\def\Rx{9.0}
% heat spreader and heat path
\fill[draftgray!14] (-0.2,4.75) rectangle (\W+0.2,5.25);
\node[lay,text=draftgray] at (\W/2,5.0) {thermal interface material and heat spreader};
\foreach \x in {1.0,2.0,3.0,10.0,11.0,12.0}{\draw[-{Stealth[length=2.2mm]},draftgray,thick] (\x,4.15) -- (\x,4.65);}
% tier 2: structural silicon | L3D | structural silicon
\filldraw[fill=draftlight,draw=draftblue] (0,3.0) rectangle (\Lx-0.1,4.05);
\filldraw[fill=draftlight,draw=draftblue] (\Rx+0.1,3.0) rectangle (\W,4.05);
\fill[pattern=north east lines,pattern color=draftblue!30] (\Lx,3.0) rectangle (\Rx,4.05);
\draw[draftblue,thick] (\Lx,3.0) rectangle (\Rx,4.05);
\node[lay] at ({(\Lx-0.1)/2},3.52) {structural silicon};
\node[lay] at ({(\Rx+0.1+\W)/2},3.52) {structural silicon};
\node[lay,fill=white,inner sep=2pt] at ({(\Lx+\Rx)/2},3.52) {L3D: 64\,MB extended L3\\(face-down)};
% bond interfaces
\draw[draftblue,line width=1.4pt] (\Lx,2.97) -- (\Rx,2.97);
\draw[draftgray,densely dashed,line width=0.9pt] (0,2.97) -- (\Lx-0.1,2.97) (\Rx+0.1,2.97) -- (\W,2.97);
% tier 1: thinned CCD, face-down
\filldraw[fill=draftlight!55,draw=draftblue,thick] (0,1.55) rectangle (\W,2.94);
\draw[draftblue!55] (\Lx,1.55) -- (\Lx,2.94) (\Rx,1.55) -- (\Rx,2.94);
\node[lay] at ({\Lx/2},2.2) {4 cores + private L2};
\node[lay] at ({(\Rx+\W)/2},2.2) {4 cores + private L2};
\node[lay] at ({(\Lx+\Rx)/2},1.85) {base L3, 32\,MB};
\foreach \x in {4.35,4.8,...,8.75}{\draw[draftblue,line width=0.7pt] (\x,2.3) -- (\x,2.94);}
% C4 bumps and substrate
\foreach \x in {0.45,1.15,...,12.7}{\fill[draftgray] (\x,1.38) circle (0.12);}
\filldraw[fill=draftgray!22,draw=draftgray] (-0.2,0.65) rectangle (\W+0.2,1.22);
\node[lay] at (\W/2,0.93) {package substrate};
% markers
\node[mk] at (0.35,3.75) {1};
\node[mk] at (\Lx+0.35,3.75) {2};
\node[mk] at (6.5,2.97) {3};
\node[mk] at (4.05,2.62) {4};
\node[mk] at (0.35,2.65) {5};
\node[mk] at (12.65,1.38) {6};
\node[mk] at (12.65,4.4) {7};
% key
\node[anchor=north,font=\scriptsize,text width=\W cm] at (\W/2,0.45) {%
\begin{tabular}{@{}p{0.47\linewidth}p{0.47\linewidth}@{}}
\textbf{1} Structural silicon dies, oxide-bonded to the CCD: support the thinned CCD and conduct heat from the cores. &
\textbf{2} L3D: 7\,nm, 41\,mm$^2$, 64\,MB of L3 over the CCD's central L3 region. \\[2pt]
\textbf{3} Hybrid bond: direct copper-to-copper connection (TSMC SoIC) between CCD and L3D. &
\textbf{4} Through-silicon vias in the thinned CCD; minimum 9\,\textmu m pitch. \\[2pt]
\textbf{5} CCD: Zen~3, 7\,nm, 81\,mm$^2$, eight cores; the shared L3 grows from 32 to 96\,MB. &
\textbf{6} C4 bumps connect the face-down CCD to the substrate. \\[2pt]
\multicolumn{2}{@{}p{0.96\linewidth}@{}}{\textbf{7} Heat from the cores rises through the structural silicon to the heat spreader.}
\end{tabular}};
\end{tikzpicture}
\caption[Cross-section of the Zen~3 3D V-Cache stack]{Schematic cross-section of the Zen~3 3D V-Cache stack, drawn from AMD's descriptions \cite[slides~10--12, 14, 16--17]{wuu2022vcache} and \cite{amd2022epyc7003}. Not to scale; layer thicknesses and the number of TSVs are illustrative.}
\label{fig:vcache-stack}
\end{figure}
